% s-184-p184-m002; original PDF 184, 185
\Aphorism{s-184-p184-m002}{274}
Das Ideal eines guten Gesprächs ist dasjenige, das wie jedes gute Ding auf einem Boden eigener Ekstase erwächst, von dem eine Verpflanzung in niedere Realitätsgattungen gar nicht möglich ist. Gedanken, die innerhalb eines solchen Gesprächs zu Worten kommen, müssen anderswo aufhören, Gedanken zu sein: ihre Wahrheit ist der befeuerte Sprechklang, in dem sie das Licht der Ideenwelt erblicken, ihr Fundament das schwebende Gleichgewicht, das sie inmitten des Rede-Luftschlosses einnehmen, ihr Gehalt die Fülle der tausend intimen Beziehungen von Ort, Zeit, Anlass, Theilnehmern des Gesprächs, ihr Glanz das Durchschimmern der tausend Nüancen augenblicklicher Stimmung. Dies alles abgerechnet, bleibt nichts, vielleicht ein logischer Fehlschluss oder ein Feuilleton: aber was bedeutet dies gegen die volle Realität des Gedankens in jenem Einen Augenblick, der zugleich Zeugung, Empfängniss und Geburt in sich schloss? — Nun, diese Gesprächs-Ekstase ist von Seiten der Alltagswirklichkeit nicht vorzubereiten, nicht stufenweise zu ersteigen, sondern sie muss von selbst aufleuchten. Ein Verständniss, das mit einer vorangehenden Einigung über theoretische Kapitel „angebahnt“ wird, ist kein Verständniss; ein Gespräch, in dem nicht ohne weiteres improvisirt werden kann, ist kein Gespräch, sondern eine monologische Auseinandersetzung mit vertheilten Rollen, ein collegium, kein colloquium. Das „lehrreiche Gespräch“ ist, wie die didaktische Poesie, eine Abart der Kunst, zur Erbauung derer, die für Kunst keine Organe haben. Vorausgesetzt wird dabei, dass zwischen den Unterrednern eine gewisse Gleichartigkeit oder Polarität von Bildung, Erfahrungsumfang, geistigem Wuchs und Eigenthum bestehe, damit sie einander mit den Früchten augenblicklicher Eingebung bewirthen können, ohne über Baum, Erdreich und Klima erklärende Anmerkungen unter den Strich zu setzen, ich meine ohne einander ihre Systeme und „Weltanschauungen“ vorzutragen. Ein gutes Gespräch — vielleicht fasst dies alles zusammen — zeitigt nur solche Gedanken, die ihrem Urheber ebenso neu sind wie seinen Zuhörern, die überhaupt nicht dem Einzelnen, sondern allen Gesprächstheilnehmern wie einem lebendigen Collectivum als geistiges Eigenthum zuzurechnen sind.
\par

